% !TeX root = ../main.tex

\section{Introducción y motivación}
Uno de los ámbitos más importantes en la astrofísica actual es el estudio de nuestro propio Sol. Los múltiples fenómenos físicos que presenta nos ayudan a comprender el comportamiento de otras estrellas, convirtiéndolo en un área de investigación muy activa. Para profundizar en su estudio, se han desarrollado observatorios como Sunrise \cite{Barthol2011}, un telescopio montado en un globo estratosférico que permite esquivar gran parte de la atmósfera terrestre. Gracias a instrumentos a bordo como TuMag \cite{SUNRISETuMagTUMAGData}, es posible obtener imágenes magnéticas solares de alta calidad y resolución espacial.

Sin embargo, el problema a resolver radica en que, a pesar de operar por encima de la atmósfera, el propio instrumento sufre los efectos de la difracción de la luz y de aberraciones ópticas. Para lograr imágenes nítidas de regiones muy pequeñas del Sol, es necesario contrarrestar estos defectos. A este proceso computacional que corrige la degradación espacial de las imágenes se le llama deconvolución.

Un aspecto fundamental del estudio solar es su campo magnético. Aunque una parte de este magnetismo se manifiesta en procesos a gran escala y muy energéticos, como las manchas solares, una fracción nada despreciable se encuentra en lo que se conoce como el ``Sol en calma'' \cite{BellotRubio2019}. En esta región destacan los gránulos solares, células convectivas de plasma situadas en la fotosfera. Estas estructuras presentan un campo magnético débil (del orden de cientos de Gauss) y son un gran foco de investigación, ya que aún se debate si su magnetismo es generado por el propio movimiento del plasma a nivel local, o si es un reflejo del campo magnético interno que emerge a la superficie.

Describir teóricamente la interacción entre el campo magnético y la luz es un reto complejo. En la literatura se han desarrollado modelos que, a través de los parámetros de Stokes, nos permiten medir estos campos magnéticos. Una de estas aproximaciones teóricas, y el pilar fundamental de este proyecto, es la aproximación de campo débil \cite{cteGeff}. Utilizando desarrollos de Taylor de primer orden, esta aproximación nos permite deducir el campo magnético de forma directa a partir de los perfiles de Stokes, facilitando el estudio del Sol en calma y de los gránulos solares.

Precisamente, afrontar el reto computacional y físico de deconvolucionar estas imágenes integrando la información del campo magnético ha sido el objetivo principal de este proyecto. Poder trabajar con datos de misiones aeroespaciales, conocer la investigación actual de primera mano, aprender de los investigadores del Departamento de Física Solar del Instituto de Astrofísica de Andalucía (IAA) y contribuir con un nuevo algoritmo han sido las principales motivaciones de este trabajo.

Para lograr este objetivo, el presente documento se estructura de la siguiente manera: en primer lugar, se explican los fundamentos teóricos sobre datos espectrales y transferencia radiativa, necesarios para derivar la expresión del campo magnético objeto de estudio. A continuación, se describe el problema de la degradación espacial de las imágenes y los retos asociados a la recolección de datos con telescopios, introduciendo así el concepto de deconvolución. Después, se presenta el estado actual de las técnicas de deconvolución más clásicas aplicadas a imágenes convencionales. Tras ello, se propone y desarrolla un nuevo método para restaurar los mapas de campo magnético basándose en la literatura científica reciente. Posteriormente, se analizan tanto los algoritmos tradicionales (implementados como parte de este trabajo) como el nuevo método propuesto, evaluando sus resultados en datos sintéticos mediante diversas métricas. Por último, se aplica el método desarrollado a datos observacionales reales recogidos por el instrumento TuMag, evaluando cualitativamente los resultados obtenidos. Además, las tecnologías y herramientas computacionales empleadas en el desarrollo de este trabajo se detallan en el Apéndice \ref{sec:tecnologias_usadas}.

\subsection{Aportación personal del trabajo}
Para aclarar el alcance de este trabajo, es importante distinguir entre los conceptos extraídos de la literatura y el trabajo desarrollado de forma autónoma. Toda la fundamentación teórica (tanto física como matemática) procede de la bibliografía citada. Por su parte, la aportación personal de este proyecto se centra en los siguientes puntos:

\begin{itemize}
    \item \textbf{Programación de todos los algoritmos:} Aunque los métodos clásicos de deconvolución no se han desarrollado teóricamente, se han programado desde cero todos los algoritmos involucrados en este trabajo, evitando el uso de ``cajas negras'' preexistentes.
    \item \textbf{Aplicación del método de van Noort:} Se ha adaptado e implementado de manera teórica el modelo directo derivado del método de van Noort (descrito en la segunda mitad de la Sección 5), aplicándolo específicamente a la recuperación del campo magnético.
    \item \textbf{Diseño de pruebas sintéticas y análisis de resultados:} Se han concebido y diseñado todos los experimentos sobre datos sintéticos para validar el correcto funcionamiento de los métodos desarrollados, simulando su respuesta frente a ópticas reales y ruido. Posteriormente, se han aplicado métricas de evaluación para analizar exhaustivamente los resultados computacionales obtenidos.
    \item \textbf{Aplicación a datos reales:} Finalmente, se ha aplicado la solución desarrollada a datos empíricos provenientes del instrumento TuMag, interpretando los resultados y redactando las conclusiones del trabajo.
\end{itemize}